\documentclass[11pt,a4paper]{article}

\usepackage[utf8]{inputenc}
\usepackage[T1]{fontenc}
\usepackage[french]{babel}
\usepackage[a4paper,margin=2cm]{geometry}
\usepackage{amsmath,amssymb}
\usepackage{lmodern}
\usepackage{fancyhdr}
\usepackage{lastpage}
\usepackage{enumitem}

\setlength{\headheight}{14pt}
\pagestyle{fancy}
\fancyhf{}
\fancyhead[L]{PCSI}
\fancyhead[C]{Devoir surveillé -- Algèbre linéaire}
\fancyhead[R]{Page \thepage/\pageref{LastPage}}
\fancyfoot[C]{Durée : 4 heures}

\begin{document}

\textbf{Consignes générales.} La qualité de la rédaction, la rigueur des raisonnements et la clarté des démonstrations seront fortement valorées.  
Toute réponse non justifiée sera peu valorisée.  
Toute initiative pertinente sera prise en compte.

\bigskip

\textbf{Répartition indicative :}

\begin{itemize}
\item Problème 1 : 40 pts
\item Problème 2 : 30 pts
\item Problème 3 : 34 pts
\end{itemize}

\bigskip
\section*{Rappels utiles}

Les rappels suivants ne dispensent pas d’une rédaction rigoureuse. Ils visent uniquement à préciser certaines notions utilisées dans le sujet.

\bigskip
\textbf{1. Codimension}

Si $F$ est un sous-espace d’un espace vectoriel de dimension finie $E$, on définit :
\[
\mathrm{codim}(F) = \dim(E) - \dim(F).
\]

\bigskip
\textbf{2. Barycentre}

Pour des vecteurs $A,B,C$, le barycentre est donné par :
\[
G = \frac{1}{3}(A + B + C).
\]

\bigskip

\textbf{3. Coplanarité}

Des points sont coplanaires si leurs vecteurs position appartiennent à un même sous-espace vectoriel de dimension au plus 2.

\bigskip

\textbf{4. Espace vectoriel de dimension infinie}

Un espace vectoriel est dit de dimension infinie s’il n’admet aucune base finie.
\bigskip

%==================================================
\section*{Problème 1 — Suites d’espaces et structure}

Pour tout $n \geq 1$, on pose $E_n = \mathbb{R}^n$.

On identifie $E_n$ à un sous-espace de $E_{n+1}$ via :
\[
(x_1,\dots,x_n) \mapsto (x_1,\dots,x_n,0)
\]

On obtient une suite croissante :
\[
E_1 \subset E_2 \subset \cdots \subset E_n \subset E_{n+1} \subset \cdots
\]

%----------------------------------
\subsection*{Partie A — Structure élémentaire (12 pts)}

\begin{enumerate}
\item Montrer que $E_n$ est un sous-espace vectoriel de $E_{n+1}$. \hfill (2 pts)
\item Déterminer $\dim(E_n)$. \hfill (1 pt)
\item Montrer que $E_n$ est de codimension 1 dans $E_{n+1}$. \hfill (5 pts)
\end{enumerate}

%----------------------------------
\subsection*{Partie B — Décomposition canonique (12 pts)}

On note $e_{n+1} = (0,\dots,0,1)$.

\begin{enumerate}
\setcounter{enumi}{3}
\item Montrer que $E_{n+1} = E_n \oplus \mathrm{Vect}(e_{n+1})$. \hfill (6 pts)
\item En déduire une construction inductive d’une base de $E_n$. \hfill (3 pts)
\item Montrer que toute base de $E_{n+1}$ contenant une base de $E_n$ s’obtient en ajoutant un vecteur hors de $E_n$. \hfill (3 pts)
\end{enumerate}

%----------------------------------
\subsection*{Partie C — Sous-espaces emboîtés (14 pts)}

Soit $(F_n)$ une suite telle que :
\[
F_n \subset E_n \quad \text{et} \quad F_n \subset F_{n+1}
\]

\begin{enumerate}
\setcounter{enumi}{6}
\item Montrer que $(\dim F_n)$ est croissante. \hfill (2 pts)
\item Montrer que $\dim(F_n) \leq n$. \hfill (1 pt)
\item On suppose que $\dim(F_{n+1}) = \dim(F_n)$. Montrer que $F_{n+1} = F_n$. \hfill (5 pts)
\item En déduire que la suite $(F_n)$ est stationnaire à partir d’un certain rang si $\dim(F_n)$ est bornée. \hfill (4 pts)
\end{enumerate}

%----------------------------------
\subsection*{Partie D — Passage à la limite (12 pts)}

On pose :
\[
E = \bigcup_{n \geq 1} E_n
\]

\begin{enumerate}
\setcounter{enumi}{10}
\item Montrer que $E$ est un espace vectoriel. \hfill (2 pts)
\item Montrer que $E$ n’admet pas de famille génératrice finie. \hfill (5 pts)
\item Montrer que toute famille libre de $E$ est finie. \hfill (3 pts)
\item Discuter la notion de dimension de $E$. \hfill (4 pts)
\end{enumerate}

\bigskip

%==================================================
\section*{Problème 2 — Familles libérées et structure}

Soit $E$ un espace vectoriel.

\textbf{Définition.} Une famille $(u_1,\dots,u_p)$ est dite \textbf{libérée} si :
\[
\forall k,\quad u_k \notin \mathrm{Vect}(u_1,\dots,u_{k-1})
\]

%----------------------------------
\subsection*{Partie A — Nature de ces familles (12 pts)}

\begin{enumerate}
\item Montrer qu’une famille libérée est libre. \hfill (3 pts)
\item La réciproque est-elle vraie ? \hfill (4 pts)
\item Donner un exemple de famille libre non libérée. \hfill (2 pts)
\end{enumerate}

%----------------------------------
\subsection*{Partie B — Réorganisation et structure (12 pts)}

\begin{enumerate}
\setcounter{enumi}{3}
\item Montrer que toute famille libre peut être réordonnée en une famille libérée. \hfill (5 pts)
\item Montrer que toute famille libérée maximale est une base. \hfill (4 pts)
\item Montrer que toute base est une famille libérée. \hfill (2 pts)
\end{enumerate}

%----------------------------------
\subsection*{Partie C — Caractérisations profondes (10 pts)}

\begin{enumerate}
\setcounter{enumi}{6}
\item Montrer qu’une famille est une base si et seulement si elle est libérée et génératrice. \hfill (6 pts)
\item Montrer qu’une famille libérée de $n$ vecteurs dans $\mathbb{R}^n$ est une base. \hfill (3 pts)
\end{enumerate}

%----------------------------------
\subsection*{Partie D — Dimension et contraintes (10 pts)}

\begin{enumerate}
\setcounter{enumi}{8}
\item Montrer que toute famille libérée de $\mathbb{R}^n$ contient au plus $n$ vecteurs. \hfill (4 pts)
\item Montrer que toute famille libérée maximale a exactement $n$ vecteurs. \hfill (2 pts)
\end{enumerate}

\newpage
%==================================================
\section*{Problème 3 — Suites de figures dans $\mathbb{R}^3$ (34 pts)}

On considère deux vecteurs $u,v$ non colinéaires de $\mathbb{R}^3$.

Pour tout entier $n \geq 1$, on définit :
\[
u_n = u + n v, \quad v_n = v + n u
\]

On associe les points :
\[
A_n = u_n, \quad B_n = v_n, \quad C_n = u_n + v_n
\]

On note $\mathcal{F}_n$ la figure formée par les points $A_n,B_n,C_n$.

%----------------------------------
\subsection*{Partie A — Étude des premières figures (8 pts)}

\begin{enumerate}
\item Calculer $\overrightarrow{A_nB_n}$ et $\overrightarrow{A_nC_n}$. \hfill (2 pts)
\item Montrer que $A_n,B_n,C_n$ ne sont pas alignés. \hfill (3 pts)
\item En déduire que $\mathcal{F}_n$ est un triangle pour tout $n$. \hfill (3 pts)
\end{enumerate}

%----------------------------------
\subsection*{Partie B — Structure vectorielle (8 pts)}

\begin{enumerate}
\setcounter{enumi}{3}
\item Exprimer $\overrightarrow{A_nB_n}$ et $\overrightarrow{A_nC_n}$ en fonction de $u$ et $v$. \hfill (3 pts)
\item Montrer que tous les triangles $\mathcal{F}_n$ sont contenus dans un même plan vectoriel. \hfill (3 pts)
\item Déterminer ce plan. \hfill (2 pts)
\end{enumerate}

%----------------------------------
\subsection*{Partie C — Évolution de la figure (8 pts)}

\begin{enumerate}
\setcounter{enumi}{6}
\item Montrer que $\overrightarrow{A_nB_n} = (n-1)u + (1-n)v$. \hfill (2 pts)
\item Étudier l’évolution de la direction de $\overrightarrow{A_nB_n}$ lorsque $n$ varie. \hfill (3 pts)
\item Existe-t-il $n$ pour lequel $A_nB_n$ est parallèle à $u-v$ ? \hfill (3 pts)
\end{enumerate}

%----------------------------------
\subsection*{Partie D — Centres et symétrie (6 pts)}

On définit :
\[
G_n = \frac{1}{3}(A_n + B_n + C_n)
\]

\begin{enumerate}
\setcounter{enumi}{9}
\item Calculer $G_n$ en fonction de $u$ et $v$. \hfill (2 pts)
\item Étudier la suite $(G_n)$. \hfill (2 pts)
\item Interpréter géométriquement le mouvement des centres. \hfill (2 pts)
\end{enumerate}

%----------------------------------
\subsection*{Partie E — Suite d’ensembles  (4 pts)}

On pose :
\[
S_n = \mathrm{Vect}(u_n,v_n)
\]

\begin{enumerate}
\setcounter{enumi}{12}
\item Montrer que $S_n$ est indépendant de $n$. \hfill (2 pts)
\item Interpréter ce résultat géométriquement. \hfill (2 pts)
\end{enumerate}

\bigskip

\textbf{Total indicatif : 104 pts}

\end{document}